%!TEX program = xelatex
\documentclass[lang=cn,a4paper,newtx]{elegantpaper}
\title{阈值控制策略分析：补充材料}
\author{XXX \and XXX}
\institute{School of Mathematics and Statistics, Shaanxi Normal University}
\date{}
\usepackage{amsmath,amssymb,array,tabularx,makecell,threeparttable,siunitx,xurl}
\renewcommand{\arraystretch}{1.25}
\emergencystretch=2em
\addbibresource{references.bib}
\renewcommand{\thetable}{S\arabic{table}}
\begin{document}
\maketitle
本文件列出基准参数分析、初值设定比较及西安阈值敏感性结果，并补充社区感染上限与医疗需求的联系和阈值取值的参考依据。数值设置、未取整结果及历史程序对账见随文复现说明。
\begin{table}[htbp]
\centering\small
\setlength{\tabcolsep}{3.5pt}
\renewcommand{\arraystretch}{1.25}
\begin{threeparttable}
\caption{基准参数下代表性控制结果}
\label{tab:sup:baseline}
\begin{tabular*}{\linewidth}{@{\extracolsep{\fill}}lS[table-format=2.4]S[table-format=3.4]S[table-format=3.4]S[table-format=1.4]S[table-format=2.4]S[table-format=3.4]@{}}
\toprule
参数取值 & {$t_1$（d）} & {$\Delta t$（d）} & {$t_{\rm end}$（d）} & {$q_{\max}$} & {$J$（d）} & {\shortstack{总累计感染\\（人）}}\\
\midrule
\multicolumn{7}{@{}l}{固定 $\eta/N=5\%$，改变 $c_0$}\\
\addlinespace[2pt]
$c_0\approx3.3470$ & 29.9176 & 2.0432 & 71.4690 & 0.0783 & 0.0053 & 382.5525\\
$c_0=4$ & 15.3697 & 6.8658 & 58.7277 & 0.3413 & 0.4374 & 360.7318\\
$c_0=5$ & 9.3039 & 7.5778 & 49.9657 & 0.4946 & 1.0561 & 342.5629\\
$c_0=8$ & 4.3073 & 6.6308 & 38.0095 & 0.6936 & 2.0006 & 303.5829\\
$c_0=10$ & 3.1754 & 5.9026 & 33.7732 & 0.7565 & 2.2236 & 285.7244\\
$c_0=12$ & 2.5150 & 5.3007 & 30.7687 & 0.7978 & 2.3110 & 271.9295\\
\addlinespace[5pt]
\multicolumn{7}{@{}l}{固定 $c_0=10$，改变 $\eta/N$}\\
\addlinespace[2pt]
$\eta/N=0.2\%$ & 0.3602 & 152.0574 & 180.5821 & 0.7734 & 60.8005 & 173.0314\\
$\eta/N=0.6\%$ & 1.3001 & 50.5672 & 86.5579 & 0.7721 & 20.1265 & 195.3349\\
$\eta/N=1\%$ & 1.7406 & 30.2685 & 64.9343 & 0.7708 & 11.9912 & 208.9166\\
$\eta/N=2\%$ & 2.3459 & 15.0431 & 46.8029 & 0.7674 & 5.8888 & 233.7810\\
\bottomrule
\end{tabular*}
\begin{tablenotes}[flushleft]\footnotesize
\item 注：表中列出基准参数下的结果；首行 $c_0=c_{0,\min}+0.02$，近似数不作重新求解的输入。第一组阈值为 $5\%$；正文多阈值曲线另取 $0.2\%,0.6\%,1\%,2\%$。
\end{tablenotes}
\end{threeparttable}
\end{table}
\clearpage
\begin{table}[htbp]
\centering\small
\setlength{\tabcolsep}{3.5pt}
\renewcommand{\arraystretch}{1.25}
\begin{threeparttable}
\caption{不同初值设定下的观测窗口比较}
\label{tab:sup:initial}
\begin{tabular*}{\linewidth}{@{\extracolsep{\fill}}lS[table-format=1.6]S[table-format=7.0]S[table-format=5.2]S[table-format=5.2]@{}}
\toprule
数据或初值设定 & {$I_0$（人）} & {\shortstack{累计病例\\（40 d）}} & {\shortstack{社区日新增峰值\\（例）}} & {\shortstack{隔离日新增峰值\\（例）}}\\
\midrule
观测数据 & {---} & 2050 & 60 & 141\\
最小二乘拟合 & 0.001007 & 2099 & 53.88 & 115.93\\
固定 $I_0=1$ & 1 & 1083553 & 36873.64 & 63966.50\\
\bottomrule
\end{tabular*}
\begin{tablenotes}[flushleft]\footnotesize
\item 注：观测窗口为 2021 年 12 月 9 日至 2022 年 1 月 17 日的 40 个日区间。模型结果采用统一的归一化方程、DOP853 求解器和分量容差；拟合行使用未取整的最小二乘初值，另一模型行固定 $I_0=1$。TDINN 拟合轨迹在 $t=40$ 的总累计为 $2098.9747$，在动态清零时刻 $45.2874$ d 为 $2105.6330$。容差收紧且最大步长减半后，清零时总累计的变化为 $3.61\times10^{-8}$ 人；该差异仅对应 TDINN 基准的清零累计量。日新增峰值不等于社区现存感染者峰值。
\end{tablenotes}
\end{threeparttable}
\end{table}
\clearpage
\begin{table}[htbp]
\centering\small
\setlength{\tabcolsep}{3.5pt}
\renewcommand{\arraystretch}{1.25}
\begin{threeparttable}
\caption{西安参数下不同阈值的控制结果}
\label{tab:sup:eta}
\begin{tabular*}{\linewidth}{@{\extracolsep{\fill}}S[table-format=1.4]S[table-format=1.4]S[table-format=4.2]S[table-format=4.2]S[table-format=3.4]S[table-format=7.0]@{}}
\toprule
{$\eta/N$} & {$q_{\max}$} & {$\Delta t$（d）} & {$t_{\rm end}$（d）} & {$J$（d）} & {总累计感染（人）}\\
\midrule
0.0001 & 0.8470 & 1710.66 & 2209.55 & 826.4066 & 2155649\\
0.0002 & 0.8469 & 855.09 & 1243.15 & 412.9141 & 2182123\\
0.0005 & 0.8466 & 341.75 & 620.77 & 164.8184 & 2234666\\
0.0010 & 0.8462 & 170.63 & 389.32 & 82.1192 & 2293939\\
0.0015 & 0.8458 & 113.59 & 304.01 & 54.5524 & 2339486\\
0.0020 & 0.8454 & 85.07 & 258.10 & 40.7686 & 2377943\\
0.0030 & 0.8445 & 56.55 & 208.39 & 26.9842 & 2442608\\
0.0040 & 0.8436 & 42.29 & 181.18 & 20.0913 & 2497306\\
0.0050 & 0.8428 & 33.73 & 163.65 & 15.9550 & 2545663\\
0.0075 & 0.8405 & 22.32 & 138.10 & 10.4383 & 2649240\\
0.0100 & 0.8382 & 16.61 & 123.89 & 7.6781 & 2737382\\
\bottomrule
\end{tabular*}
\begin{tablenotes}[flushleft]\footnotesize
\item 注：数值采用与正文一致的西安参数，按仅隔离控制的解析表达及数值求积计算；$q_{\max}=q_c(t_1)$。$\eta/N$ 为比例而非百分数，0.002对应0.2\%。阈值网格用于比较不同设计情景。累计量计算至各策略的动态清零终点。长时间结果沿用固定参数模型的外推解释，不作为现实预测。
\end{tablenotes}
\end{threeparttable}
\end{table}

\clearpage
\renewcommand{\thesection}{S\arabic{section}}
\section{社区感染上限与医疗需求的联系}
\label{sec:sup:medical-link}
本节在正文传播模型之外补充一个医疗需求关系，说明如何在附加假设下根据容量选取社区感染上限。

两类感染仓室的新增流入之和为
\[
\frac{\beta c(t)(1-q(t))S(t)I(t)}{N}
+\frac{\beta c(t)q(t)S(t)I(t)}{N}
=\frac{\beta c(t)S(t)I(t)}{N}.
\]
若全过程满足 $0\le c(t)\le c_0$ 和 $I(t)\le\eta$，由无回流模型中的 $S(t)\le S_0$，得到
\[
0\le\frac{\beta c(t)S(t)I(t)}{N}
\le\frac{\beta c_0S_0}{N}\eta.
\]
这一上界包含进入 $I$ 和 $I_q$ 的全部新增感染。

设非负可测函数 $h(a)$ 表示一例新感染者在感染后第 $a$ 天所需的平均 ICU 床位数，并假定它对两类病例、不同感染时点及所比较的干预分配共同适用。令
\[
\int_0^\infty h(a)\,da=pL<\infty,
\]
其中 $p\in(0,1]$ 为需要 ICU 的概率，$L>0$ 为需要 ICU 者累计占用时长的均值；$h$ 可包含感染至入院的延迟。根据既往流入与停留过程计算床位占用的建模方法可参见文献\cite{Baas2021Occupancy}。

初始时刻之前已经感染者此后的平均需求记为 $D_0(t)$，并假定各策略均满足 $0\le D_0(t)\le D_{\mathrm{init}}<\infty$。该项包括初始病例以后才发生的入院和后续占用，也包括已从传播模型移出但仍需救治的病例。未受容量拒收限制的平均需求为
\[
D(t)=D_0(t)+\int_0^t h(t-u)
\frac{\beta c(u)S(u)I(u)}{N}\,du.
\]
利用非负性及总流入上界，有
\[
\begin{aligned}
D(t)
&\le D_{\mathrm{init}}+
\frac{\beta c_0S_0}{N}\eta\int_0^t h(a)\,da\\
&\le D_{\mathrm{init}}+pL\frac{\beta c_0S_0}{N}\eta,
\qquad t\ge0.
\end{aligned}
\]
设 $C>D_{\mathrm{init}}$ 为全过程固定的本病可用容量，已扣除其他病种占用和安全预留，本病初始病例的未来负担由 $D_{\mathrm{init}}$ 单独计入。于是
\[
\eta\le\frac{C-D_{\mathrm{init}}}{pL\,\beta c_0S_0/N}
\]
是平均需求满足 $D(t)\le C$ 的充分条件。该界对允许的干预分配共同成立，可能偏保守；它针对上述平均需求关系，不给出随机超载概率。按此条件选取阈值后，仍须检验正文中的触发条件和干预能力条件。

\section{代表性阈值的量级参考}
\label{sec:sup:threshold-scale}
正文西安算例采用 $N=13{,}163{,}000$\cite{XianStat2021}、$\eta=0.002N=26{,}326$。本节保留其比例情景的参考依据，并与上一节的医疗需求充分条件区分。

陈胤孜等\cite{Chen2021ICUForecast}基于既往资料预测，2021 年全国每10万人综合 ICU 床位约为 $4.37$ 张；这一数值为全国预测量，而非西安同期可用容量的实测值。Angulo 等\cite[补充材料 S4.1]{Angulo2021}对早期 COVID-19 采用需要重症监护的感染者比例名义值 $0.60\times0.15\times0.28=0.0252$。将前者作为人均容量参考、不扣除既有占用，并将后者移作参考系数 $\rho_{\mathrm{ICU}}$，按比例情景
\[
\eta=\frac{C}{\rho_{\mathrm{ICU}}},\qquad
\theta=\frac{C/N}{\rho_{\mathrm{ICU}}}
\]
得到
\[
\theta=\frac{4.37/10^5}{0.0252}\approx1.73\times10^{-3},
\qquad \eta\approx2.28\times10^4.
\]
正文在这一量级附近取 $\theta=0.002$ 作代表性设计情景，而非对上式作保守取整。按相同参考系数反算，它对应每10万人 $5.04$ 张、全市约 $663$ 张的名义容量。

上述比例移植用于说明阈值量级，不是当地医疗容量的临床校准。以确诊病例或全部感染者为分母的 ICU 入院风险可能不同\cite{Pung2023Undetected}；此处的参考系数未依据共同需求函数和初始负担估计，也不等同于第\ref{sec:sup:medical-link}节的 $pL\beta c_0S_0/N$。两节分别给出一个历史比例情景和一个附加医疗假设下的充分条件，现有算例采用前者附近的设计阈值开展控制比较。

\printbibliography[title={参考文献}]
\end{document}
